\documentclass[11pt,a4paper]{ctexart}
\usepackage[a4paper,margin=1.78cm,headheight=15pt]{geometry}
\usepackage{amsmath,amssymb,mathtools,bm}
\usepackage{booktabs,longtable,tabularx,array}
\usepackage{enumitem}
\usepackage[most]{tcolorbox}
\usepackage{tikz}
\usetikzlibrary{arrows.meta,positioning,fit,calc,decorations.pathreplacing}
\usepackage{xcolor}
\usepackage{fancyhdr}
\usepackage{hyperref}
\usepackage{bookmark}
\usepackage{microtype}

\newcolumntype{Y}{>{\raggedright\arraybackslash}X}
\newcolumntype{L}[1]{>{\raggedright\arraybackslash}p{#1}}
\setlength{\parindent}{2em}
\setlength{\parskip}{0.34em}
\setlength{\emergencystretch}{3em}
\renewcommand{\arraystretch}{1.22}
\setlength{\tabcolsep}{3pt}
\setlist[itemize]{itemsep=0.18em,topsep=0.35em,leftmargin=2em}
\setlist[enumerate]{itemsep=0.18em,topsep=0.35em,leftmargin=2em}

\definecolor{deep}{RGB}{42,58,73}
\definecolor{soft}{RGB}{245,247,249}
\definecolor{accent}{RGB}{104,76,55}
\definecolor{greenish}{RGB}{57,92,76}
\definecolor{warn}{RGB}{136,68,63}
\definecolor{purple}{RGB}{87,72,116}
% 训练图与 Markdown SVG 共用的语义色；本册 PDF 固定使用亮色映射。
\definecolor{themebg}{HTML}{FAFAF7}
\definecolor{themefg}{HTML}{111111}
\definecolor{themegray}{HTML}{666666}
\definecolor{themecurve}{HTML}{1D63B8}
\definecolor{themealert}{HTML}{C53030}
\definecolor{themeamber}{HTML}{B86E00}
\hypersetup{colorlinks=true,linkcolor=deep,urlcolor=accent,pdftitle={极限与连续：邻域、尺度与存在性}}
\pagestyle{fancy}
\fancyhf{}
\lhead{\small\color{deep}\textbf{数学一 · 高等数学专题02}}
\rhead{\small\color{deep}邻域 · 尺度 · 存在性}
\cfoot{\small\thepage}
\renewcommand{\headrulewidth}{0.4pt}
\newtcolorbox{corebox}[1]{breakable,colback=soft,colframe=deep,title=\textbf{#1},boxrule=0.8pt,arc=2mm}
\newtcolorbox{methodbox}[1]{breakable,colback=white,colframe=greenish,title=\textbf{#1},boxrule=0.7pt,arc=1.5mm}
\newtcolorbox{warnbox}[1]{breakable,colback=white,colframe=warn,title=\textbf{#1},boxrule=0.7pt,arc=1.5mm}
\newtcolorbox{examplebox}[1]{breakable,colback=white,colframe=purple,title=\textbf{#1},boxrule=0.7pt,arc=1.5mm}
\newcommand{\kw}[1]{\textbf{\color{deep}#1}}

\begin{document}
\begin{center}
{\LARGE\textbf{极限与连续：邻域、尺度与存在性}}\\[0.4em]
{\large 从“代入算值”到“控制全部允许的趋近方式”}
\end{center}
\vspace{0.5em}

\begin{quote}
本册只追问一个根本问题：\textbf{当输入按题目允许的方式不断接近时，能否用同一个输出目标统一控制全部足够靠后的输入？} 极限回答“邻域是否稳定”，尺度回答“哪一部分决定稳定速度”，连续回答“邻域趋势能否与点值接上”。
\end{quote}

\begin{center}\rule{0.5\linewidth}{0.5pt}\end{center}

% ============================================================
\setcounter{section}{-1}
\section{本册定位}
% ============================================================

本册是高等数学学科总图下的专题02。专题01先确认函数对象、定义域和表示；本册随后建立极限这套共同底座，使微分、积分、级数和微分方程中的“无限逼近”拥有可检查的合法性语言。

\subsection{本册拥有}

本册独立负责：
\begin{itemize}
  \item 数列极限、函数极限、单侧极限、无穷远极限与无穷大记号；
  \item 极限存在性的邻域控制、夹逼、归结原则与反例构造；
  \item 无穷小的阶、同阶、等价、小 $o$ 与主导尺度；
  \item 七类未定式的判别、标准化与“需要哪一级信息”的计算方法地图；
  \item 连续、单侧连续、连续运算的资格条件与间断点分类；
  \item 一阶递推数列中“合法轨道 $\to$ 收敛发动机 $\to$ 不动点识别”的基本机制；
  \item 极限问题的统一检查协议与独立验证方式。
\end{itemize}

\subsection{调用与停止边界}

本册调用专题01的定义域、复合接口与表示边界，但不重复定义函数对象。下列内容只保留接口：
\begin{itemize}
  \item 有限泰勒模型、导数定义和局部线性化由专题03拥有；
  \item 洛必达的中值定理基础、闭区间存在性、零点与函数形状由专题04拥有；
  \item 变上限积分、黎曼和与反常积分由专题05拥有；
  \item 数项级数研究部分和数列的尾部，完整判敛机制由专题10拥有；
  \item 函数列的一致收敛、逐项操作等只作扩展，不占据数学一专题02主干。
\end{itemize}

工具的本体机制由对应知识归属专题负责；单一问题族中“看到什么信号、第一步做什么、什么时候换路或停止”由同目录训练 Markdown 承担。只有同一控制动作跨多个训练专题反复成立并经证据验证，才晋升《高等数学做题规则》。但本册必须拥有极限计算的\textbf{知识层方法地图}：为什么一个未定式需要标准化、不同变形保留或丢失什么信息、当前目标需要哪一级局部精度，以及为什么此时应该调用等价、夹逼、洛必达、泰勒、导数定义或累积接口。否则专题02会退化成“只讲极限定义，却把极限怎样计算全部丢给训练层”的空心模型。

% ============================================================
\section{根本问题：一个不能直接到达的状态，怎样被精确定义？}
% ============================================================

把 $x=a$ 代入 $f(x)$，得到的是\kw{点值}；让 $x$ 接近但不等于 $a$，研究的是\kw{邻域行为}。两者只有在连续时才会接上。

极限之所以出现，是因为很多目标对象不能靠一次有限代入得到：
\begin{itemize}
  \item 数列只有一项项向后走，永远没有“第无穷项”；
  \item 函数可能在 $a$ 点没有定义，但去心邻域仍呈现稳定趋势；
  \item $0/0$ 不是数值，而是两个量同时缩小时的相对尺度尚未暴露；
  \item 振荡函数可能始终有界，却没有唯一稳定输出。
\end{itemize}

朴素方案是“取几个很近的点看看”。它可以猜候选值，却不能证明\textbf{所有足够近的合法输入}都被控制。极限机制因此必须从抽样观察升级为全称约束：

\[
\boxed{\text{任给输出误差}\ \varepsilon
\ \longrightarrow\ \text{找到输入阈值}
\ \longrightarrow\ \text{阈值以后全部合法输入都满足要求}}
\]

\begin{corebox}{极限不是“越来越近”的口号，而是一份控制契约}
挑战者先给出任意精度 $\varepsilon>0$；证明者必须据此给出一个阈值 $\delta$ 或 $N$；随后挑战者可以选择阈值内的\textbf{任意}合法输入。若输出误差始终小于 $\varepsilon$，候选值才是极限。

量词顺序不能交换：阈值可以依赖 $\varepsilon$，但不能在看见挑战者选定的某个 $x$ 或 $n$ 后再临时改变。
\end{corebox}

\begin{center}
\begin{tikzpicture}[color=themefg, text=themefg]
\begin{scope}[x=2.25cm,y=2.05cm]
  % 语义参数只用于读图：a=2, delta=0.8, L=1.5, epsilon=0.55。
  % 输入 δ 带与输出 ε 带。大面积填充低不透明度，避免压住主曲线。
  \fill[themecurve, opacity=0.08] (1.20,0.18) rectangle (2.80,3.05);
  \fill[themecurve, opacity=0.08] (0.02,0.95) rectangle (4.55,2.05);

  % 坐标轴
  \draw[themefg, line width=0.9pt, -{Stealth[length=2.0mm]}]
    (0,0.18) -- (4.65,0.18) node[below left=1pt] {$x$};
  \draw[themefg, line width=0.9pt, -{Stealth[length=2.0mm]}]
    (0,0.18) -- (0,3.13) node[below left=1pt] {$y$};

  % ε / δ 边界虚线
  \draw[themegray, densely dashed, line width=0.95pt] (1.20,0.18) -- (1.20,3.02);
  \draw[themegray, densely dashed, line width=0.95pt] (2.80,0.18) -- (2.80,3.02);
  \draw[themegray, densely dashed, line width=0.95pt] (0.02,0.95) -- (4.48,0.95);
  \draw[themegray, densely dashed, line width=0.95pt] (0.02,2.05) -- (4.48,2.05);
  \draw[themegray, densely dotted, line width=0.8pt] (0.02,1.50) -- (4.48,1.50);

  % 一般化非对称曲线
  \draw[themecurve, line width=2.1pt]
    plot[domain=0.18:1.965, samples=150]
      (\x,{1.50+0.38*(\x-2)+0.18*(\x-2)*(\x-2)-0.06*(\x-2)*(\x-2)*(\x-2)});
  \draw[themecurve, line width=2.1pt]
    plot[domain=2.035:4.35, samples=170]
      (\x,{1.50+0.38*(\x-2)+0.18*(\x-2)*(\x-2)-0.06*(\x-2)*(\x-2)*(\x-2)});

  % 双向挤压逼近箭头：展示 x 向 a 靠拢的动态
  \draw[themecurve, line width=1.1pt, -{Stealth[length=1.8mm]}] (1.35,0.40) -- (1.82,0.40);
  \draw[themecurve, line width=1.1pt, -{Stealth[length=1.8mm]}] (2.65,0.40) -- (2.18,0.40);

  % x=a 处的极限值与点值分离
  \draw[themefg, fill=themebg, line width=1.0pt] (2.00,1.50) circle (2.6pt);
  \fill[themealert] (2.00,2.66) circle (2.5pt);
  \draw[themegray, densely dotted, line width=0.8pt] (0.10,2.66) -- (2.00,2.66);
  \node[text=themealert, fill=themebg, inner sep=0.8pt, anchor=west] at (2.08,2.66) {\small $f(a)$};
  \node[text=themecurve, fill=themebg, inner sep=0.8pt, anchor=south west] at (3.42,2.43) {\small $y=f(x)$};

  % 轴标签微遮罩
  \node[text=themefg, fill=themebg, inner sep=0.8pt, anchor=east] at (-0.05,2.05) {\footnotesize $L+\varepsilon$};
  \node[text=themefg, fill=themebg, inner sep=0.8pt, anchor=east] at (-0.05,1.50) {\footnotesize $L$};
  \node[text=themefg, fill=themebg, inner sep=0.8pt, anchor=east] at (-0.05,0.95) {\footnotesize $L-\varepsilon$};

  \node[text=themefg, fill=themebg, inner sep=0.8pt, anchor=north] at (1.20,0.14) {\footnotesize $a-\delta$};
  \node[text=themefg, fill=themebg, inner sep=0.8pt, anchor=north] at (2.00,0.14) {\footnotesize $a$};
  \node[text=themefg, fill=themebg, inner sep=0.8pt, anchor=north] at (2.80,0.14) {\footnotesize $a+\delta$};

  % 输入带与输出带花括号
  \draw[themegray, decorate, decoration={brace, amplitude=4pt, mirror}]
    (1.20,-0.10) -- (2.80,-0.10);
  \node[text=themecurve, fill=themebg, inner sep=0.8pt, anchor=north] at (2.00,-0.18)
    {\footnotesize 输入去心邻域：$0<|x-a|<\delta$};

  \draw[themegray, decorate, decoration={brace, amplitude=4pt}]
    (4.62,0.95) -- (4.62,2.05);
  \node[text=themecurve, fill=themebg, inner sep=0.8pt, anchor=west] at (4.72,1.50)
    {\footnotesize 输出控制：$|f(x)-L|<\varepsilon$};

  % 结论徽标
  \node[
    draw=themegray,
    rounded corners=1.5mm,
    fill=themebg,
    inner xsep=7pt,
    inner ysep=4pt,
    anchor=north,
    text=themefg
  ] at (2.30,-0.48)
    {\small $0<|x-a|<\delta\implies|f(x)-L|<\varepsilon$};
\end{scope}
\end{tikzpicture}

\end{center}

这张母图把量词契约翻译成几何动作：先由任意给定的 $\varepsilon$ 锁定输出水平带，再寻找一个足够小的 $\delta$，使整个去心输入邻域都被映入该输出带。图中特意让 $f(a)\ne L$，用来阻断“极限必须等于点值”的错误直觉。

% ============================================================
\section{母模型：趋近规范、统一控制与点值接缝}
% ============================================================

本册把极限与连续压缩为五个依次发生的判断：
\[
\boxed{
\text{趋近规范}
\to \text{合法输入}
\to \text{输出控制}
\to \text{路径一致}
\to \text{点值接缝}
}
\]

这里的箭头是\textbf{检查顺序}：
\begin{enumerate}
  \item \kw{趋近规范}：明确 $n\to\infty$、$x\to a^\pm$、$x\to a$ 还是 $x\to\infty$；
  \item \kw{合法输入}：由定义域与趋近方向确定真正允许的尾部或邻域；
  \item \kw{输出控制}：用界、恒等变形或主导尺度把输出压入目标邻域；
  \item \kw{路径一致}：所有允许路径是否趋向同一值；
  \item \kw{点值接缝}：若还问连续，检查这个共同邻域值是否等于点值。
\end{enumerate}

\begin{center}
\begin{tikzpicture}[
  node distance=0.8cm and 0.8cm,
  box/.style={draw=deep,rounded corners=1.8mm,minimum width=3.3cm,minimum height=1.05cm,align=center,fill=soft,font=\small,inner sep=4pt},
  succbox/.style={draw=greenish,rounded corners=1.8mm,minimum width=3.3cm,minimum height=1.05cm,align=center,fill=white,font=\small,inner sep=4pt,thick},
  failbox/.style={draw=warn,rounded corners=1.8mm,minimum width=3.3cm,minimum height=1.05cm,align=center,fill=white,font=\small,inner sep=4pt,thick},
  gatebox/.style={draw=accent,rounded corners=1.8mm,minimum width=3.3cm,minimum height=1.05cm,align=center,fill=white,font=\small,inner sep=4pt,thick},
  arr/.style={-{Latex[length=2.2mm,width=1.5mm]},thick,draw=deep},
  warnarr/.style={-{Latex[length=2.2mm,width=1.5mm]},thick,draw=warn,densely dashed},
  succarr/.style={-{Latex[length=2.2mm,width=1.5mm]},thick,draw=greenish}
]
% 第一排：主干认知链（趋近规范 -> 合法输入 -> 输出控制 -> 路径一致）
\node[box] (app)  at (0,0)    {\textbf{1. 趋近规范}\\[2pt]\scriptsize $n\to\infty,\, x\to a^\pm,\, x\to a$};
\node[box] (dom)  at (4.2,0)  {\textbf{2. 合法输入}\\[2pt]\scriptsize 定义域截断 / 去心邻域};
\node[box] (ctrl) at (8.4,0)  {\textbf{3. 输出控制}\\[2pt]\scriptsize 主导尺度 / 夹逼 $\lvert f-L\rvert<\varepsilon$};
\node[box] (agr)  at (12.6,0) {\textbf{4. 路径一致}\\[2pt]\scriptsize 全称检验：左右/子列同值};

\draw[arr] (app) -- (dom);
\draw[arr] (dom) -- (ctrl);
\draw[arr] (ctrl) -- (agr);

% 第二排左侧：快速否定分支 (Fast Fail)
\node[failbox] (fail) at (4.2,-2.1) {\textbf{\color{warn}否定极限 (不存在)}\\[2pt]\scriptsize 发现两列/左右冲突或无界振荡};
\draw[warnarr] (dom) -- node[fill=white,inner sep=2pt,font=\scriptsize\color{warn},align=center]{构造一对冲突路径\\(存在反例即结束)} (fail);

% 第二排右侧：通过检验 -> 极限存在
\node[gatebox] (lim) at (12.6,-2.1) {\textbf{极限存在：$\lim f = L$}\\[2pt]\scriptsize 唯一确定目标状态};
\draw[succarr] (agr) -- node[fill=white,inner sep=2pt,font=\scriptsize\color{greenish}]{全部路径一致} (lim);

% 第三排：连续接缝与间断分类
\node[box] (stitch) at (12.6,-4.2) {\textbf{5. 连续接缝}\\[2pt]\scriptsize 核对点值：$L \overset{?}{=} f(a)$};
\draw[arr] (lim) -- (stitch);

\node[succbox] (cont)  at (6.3,-4.2) {\textbf{\color{greenish}连续}\\[2pt]\scriptsize 极限存在且 $L=f(a)$};
\node[failbox] (remov) at (12.6,-6.1) {\textbf{\color{warn}可去间断}\\[2pt]\scriptsize 极限存在但 $L\ne f(a)$ 或无定义};

\draw[succarr] (stitch) -- node[above=2pt,font=\scriptsize\color{greenish},fill=white,inner sep=2pt]{$L=f(a)$ (点值缝合)} (cont);
\draw[warnarr] (stitch) -- node[fill=white,inner sep=2pt,font=\scriptsize\color{warn}]{$L\ne f(a)$ (缝合失败)} (remov);

\end{tikzpicture}
\end{center}

这张图完整呈现了极限的控制与验证拓扑：
\begin{itemize}
  \item \textbf{主干控制}：数列与函数极限先确定趋近规范与合法输入；夹逼与尺度完成输出控制；左右极限和海涅定理负责路径一致；连续与间断负责点值接缝。
  \item \textbf{第一类间断}：左右极限均存在（$\lim_{x\to a^-}=\lim_{x\to a^+}\ne f(a)$ 为\textbf{可去间断}；$\lim_{x\to a^-}\ne\lim_{x\to a^+}$ 为\textbf{跳跃间断}）。
  \item \textbf{第二类间断}：至少一侧极限不存在（无穷间断、振荡间断）。
  \item \textbf{控制不对称性}：证明存在需控制全部输入；否定极限只需构造一对冲突路径或一个振荡反例。
\end{itemize}

% ============================================================
\section{同一控制契约的四种对象}
% ============================================================

\subsection{数列极限：控制离散尾部}

\[
\lim_{n\to\infty}a_n=L
\]
的严格含义是
\[
\forall\varepsilon>0,\ \exists N\in\mathbb N,\ \forall n>N,
\qquad \lvert a_n-L\rvert<\varepsilon.
\]

$N$ 之前有多少异常项都不影响极限；极限只拥有\kw{尾部状态}。这也是“修改有限项不改变数列极限”的生成原因。

\subsection{数列极限的唯一性：尾部不能同时稳定到两个分离目标}

若同一数列既满足 $a_n\to A$ 又满足 $a_n\to B$，则必有 $A=B$。反证地，若 $A\ne B$，取
\[
\varepsilon=\frac{\lvert A-B\rvert}{3}.
\]
充分靠后时同时有
\[
\lvert a_n-A\rvert<\varepsilon,
\qquad
\lvert a_n-B\rvert<\varepsilon.
\]
于是三角不等式给出
\[
\lvert A-B\rvert
\le \lvert A-a_n\rvert+\lvert a_n-B\rvert
<2\varepsilon
=\frac23\lvert A-B\rvert,
\]
矛盾。

这条定理的意义不是“极限答案只能写一个”，而是：\kw{一旦尾部真的被压进任意小的同一邻域，目标状态本身就被唯一锁定。} 因此后面所有“由不同路线得到候选极限”的题，都必须最终回到同一个目标；若两条合法子列给出不同候选，原数列就不可能收敛。

\subsection{子列：删除项不会创造新的尾部极限}

从原下标中取严格递增的整数列
\[
n_1<n_2<\cdots,
\]
得到 $\{a_{n_k}\}$，称为 $\{a_n\}$ 的一个子列。子列允许\textbf{删项}，但不允许改变原有次序。因为 $n_k\ge k$，所以 $k\to\infty$ 时必有 $n_k\to\infty$。

若
\[
a_n\to L,
\]
则任意子列都满足
\[
a_{n_k}\to L.
\]
证明只需直接调用尾部定义：原列从某个 $N$ 开始全部落入 $L$ 的任意给定邻域，而子列最终也只能从这些已经受控的尾部项中继续取值。

反过来，\textbf{若原列存在两条子列趋向不同极限，则原列必不收敛。} 这就是用奇偶子列、相位子列否定极限的理论基础。

\begin{corebox}{反向关键不是“看了几条子列”，而是“是否覆盖全部尾部”}
若只知道若干子列都趋于 $L$，一般不能推出原列趋于 $L$；因为未被这些子列观察到的下标仍可能藏着另一种行为。

真正可逆的情形是：存在有限组下标集合，它们在充分靠后时\textbf{覆盖原列的全部下标}，且每一组按原顺序形成的子列都趋于同一个 $L$。此时对任意 $\varepsilon>0$，分别给每组取一个足够靠后的阈值，再取有限个阈值的最大值，就能让原列所有足够靠后的项都落入 $L$ 的同一邻域，因此 $a_n\to L$。

最常见的特例是奇偶分解：
\[
a_{2n}\to L,
\qquad
 a_{2n+1}\to L
\quad\Longrightarrow\quad
 a_n\to L,
\]
因为奇数下标和偶数下标覆盖了全部整数。

但
\[
a_{3n}\to L,
\qquad
 a_{3n+1}\to L
\]
并不足够，因为 $3n+2$ 这一整类下标尚未被控制。最小反例可取
\[
a_{3n}=0,
\qquad
 a_{3n+1}=0,
\qquad
 a_{3n+2}=1.
\]
前两条子列都趋于 $0$，原列却不收敛。
\end{corebox}

这可以压成一个比“奇偶子列技巧”更稳定的模型：
\[
\boxed{\text{原列}\to\text{任意子列：自动继承；}\qquad
\text{子列族}\to\text{原列：必须覆盖全部尾部}.}
\]

\begin{corebox}{真题母例：2015 年把“继承”与“覆盖”放进同一道选择题}
若 $x_n\to a$，则 $x_{2n}$、$x_{2n+1}$、$x_{3n}$、$x_{3n+1}$ 都只是原列的子列，因此全部继承极限 $a$。反向若
\[
x_{2n}\to a,\qquad x_{2n+1}\to a,
\]
奇偶下标覆盖全部尾部，所以原列收敛到 $a$；但若只给
\[
x_{3n}\to a,\qquad x_{3n+1}\to a,
\]
则 $3n+2$ 仍是一个完全自由的尾部通道，无法推出原列收敛。

这道题值得作为母例保留，因为四个选项并不是四条口诀，而是在反复测试同一件事：\textbf{信息是从整体向抽样传，还是从抽样向整体拼回；若是后者，覆盖是否完整？}
\end{corebox}

\subsection{单调有界准则：把尾部自由度压成单侧逼近}

若 $\{a_n\}$ 单调递增且有上界，则它必收敛；单调递减且有下界时同理。以递增情形为例，令
\[
L=\sup\{a_n:n\in\mathbb N\}.
\]
对任意 $\varepsilon>0$，由上确界定义，存在某个 $N$ 使
\[
L-\varepsilon<a_N\le L.
\]
又因数列递增，所以对所有 $n\ge N$，
\[
L-\varepsilon<a_N\le a_n\le L,
\]
于是
\[
\lvert a_n-L\rvert<\varepsilon.
\]

这里真正消失的是\kw{尾部自由度}：单调性保证后续不能回头，有界性保证轨道不能逃向无穷；两者合在一起，才把尾部压到一个唯一终点。只给其中一项都不够：$a_n=n$ 单调但无界，$a_n=(-1)^n$ 有界但振荡。

\subsection{差趋零：只负责让两条轨道合流，不负责制造收敛}

若
\[
a_n-b_n\to0,
\]
它表达的是两条轨道的\kw{渐近距离消失}。因此只要其中一列已经有极限，就能把存在性无损传给另一列：若 $a_n\to L\in\mathbb R$，则
\[
b_n=a_n-(a_n-b_n)\to L.
\]
同样地，若 $a_n\to+\infty$ 或 $a_n\to-\infty$，加上一个趋零扰动也不会改变其无穷方向。

但差趋零本身\textbf{不会制造极限存在性}。例如
\[
a_n=b_n=(-1)^n
\]
满足 $a_n-b_n\equiv0$，两列却都不收敛。

有一个高价值的加强结构：若 $\{a_n\}$ 单调递减、$\{b_n\}$ 单调递增，并且
\[
a_n-b_n\to0,
\]
则两列一定收敛到同一个有限极限。因为 $d_n=a_n-b_n$ 单调递减且趋于 $0$，所以必有 $d_n\ge0$，即 $b_n\le a_n$。于是
\[
b_1\le b_n\le a_n\le a_1,
\]
两列分别单调有界，从而各自收敛；最后再由差趋零迫使两个极限相等。

这类题最稳定的顺序不是先猜公共极限，而是
\[
\boxed{\text{方向}\to\text{夹住}\to\text{各自存在}\to\text{差距消失}\to\text{公共极限}.}
\]

\subsection{函数在有限点的极限：控制去心邻域}

\[
\lim_{x\to a}f(x)=L
\]
表示
\[
\forall\varepsilon>0,\ \exists\delta>0,\quad
0<\lvert x-a\rvert<\delta\Longrightarrow \lvert f(x)-L\rvert<\varepsilon.
\]

条件 $0<\lvert x-a\rvert$ 明确排除了点 $a$ 本身，因此 $f(a)$ 未定义、定义成别的值，均不直接改变这个极限。

\subsection{单侧极限：允许集合被定义域截断}

右极限只允许 $a<x<a+\delta$，左极限只允许 $a-\delta<x<a$。双侧极限存在当且仅当
\[
\lim_{x\to a^-}f(x)=\lim_{x\to a^+}f(x)=L.
\]

分段点、绝对值规则切换点、取整点、分母变号点和定义域端点，都会改变合法输入集合。端点处若定义域只提供一侧，就只讨论相应单侧连续，不凭空要求另一侧。

\subsection{无穷远极限与无穷大：把阈值换成远端尾部}

\[
\lim_{x\to+\infty}f(x)=L
\]
表示对任意 $\varepsilon>0$，存在 $M$，使 $x>M$ 时都有 $\lvert f(x)-L\rvert<\varepsilon$。

而 $f(x)\to+\infty$ 表示：对任意高度 $A>0$，存在 $M$，使 $x>M$ 时都有 $f(x)>A$。它描述的是一种扩展发散状态，不是某个有限实数极限。

\begin{warnbox}{无界不等于趋于正无穷}
$f(x)=x\sin x$ 在 $x\to+\infty$ 时无界，但它反复取很大的正值和负值，不能满足“所有足够大的 $x$ 都高于任意给定阈值”。无界只否定有限范围；趋于 $+\infty$ 还要求最终保持方向一致。
\end{warnbox}

\subsection{有限极限怎样传递有界、符号与大小关系？}

有限极限不仅给出一个数值 $L$，还给出\kw{最终邻域中的结构控制}。设
\[
\lim_{x\to a}f(x)=L\in\mathbb R.
\]

首先，有限极限自动产生局部有界性。取 $\varepsilon=1$，则存在去心邻域使
\[
\lvert f(x)-L\rvert<1,
\]
于是
\[
\lvert f(x)\rvert\le \lvert L\rvert+1.
\]
所以“有有限极限 $\Rightarrow$ 局部有界”；反向当然不成立，有界振荡就是最小反例。

其次，\kw{严格离开零点的极限会保号}。若 $L>0$，取 $\varepsilon=L/2$，则充分靠近 $a$ 时
\[
f(x)>L-\frac L2=\frac L2>0.
\]
若 $L<0$ 同理最终恒负。因此极限保号真正依赖的是“极限值与 $0$ 之间存在正距离”，而不是只知道 $L\ge0$。例如 $f(x)=-x^2\to0$，却在去心邻域内始终为负。

\begin{warnbox}{先过存在性与共同定义域门，再谈次序传递}
“邻域里谁大”与“极限谁大”不是同一层信息。若命题的结论中出现 $\lim f$、$\lim g$，第一步必须确认这些极限已经由题设或前面的证明保证存在；邻域内的大小关系本身不会制造收敛性。例如 $x\to0$ 时
\[
f(x)=2+\sin\frac1x>0=g(x),
\]
但 $\lim_{x\to0}f(x)$ 不存在，因此不能仅凭 $f>g$ 写出两个极限的大小关系。

反向从极限值推邻域函数值时，在本册采用的函数极限约定下，“$\lim_{x\to a}f(x)$ 存在”已经包含 $f$ 在某个去心邻域内有定义；$g$ 同理。取两个去心邻域的交集即可得到共同可比较区域，$f(a)$、$g(a)$ 本身是否定义不影响这一步。
\end{warnbox}

\begin{corebox}{先抓住两个对比：严格号在两个方向上的命运不同}
先不要急着记“三道门”，先记住下面这两个成对结论：

\textbf{对比一：从极限值到邻域函数值，严格间隙可以保住严格号。}
\[
\boxed{\lim f=A,\ \lim g=B,\ A>B}
\quad\Longrightarrow\quad
\boxed{\text{充分靠近时 }f(x)>g(x)}.
\]
原因是 $A-B>0$ 给出了一个\kw{不会消失的正间隙}。令 $h=f-g$，则
\[
h(x)\to A-B>0,
\]
由极限保号性，充分靠近时 $h(x)>0$，也就是 $f(x)>g(x)$。

例如若
\[
f(x)\to3,\qquad g(x)\to1,
\]
那么两极限之间相差 $2$。函数值即使在附近有小幅波动，也不可能无限靠近后仍跨过这条固定间隙；最终必有 $f(x)>g(x)$。

\textbf{对比二：从邻域函数值到极限值，严格号可能在取极限时压成等号。}
\[
\boxed{\text{充分靠近时 }f(x)>g(x),\ \lim f,\lim g\text{ 都存在}}
\quad\Longrightarrow\quad
\boxed{\lim f\ge\lim g},
\]
但一般\textbf{不能}推出 $\lim f>\lim g$。

最小反例是
\[
f(x)=x^2,\qquad g(x)=0,\qquad x\to0.
\]
任意 $x\ne0$ 都有 $x^2>0$，所以去心邻域里始终 $f(x)>g(x)$；但
\[
\lim_{x\to0}f(x)=0=\lim_{x\to0}g(x).
\]
也就是说，两个函数可以在每一个靠近点都严格分开，却让这段差距不断缩小并最终在极限处压到同一个值。

因此可以先用一句话记住：
\[
\boxed{\text{极限有固定正间隙}\Rightarrow\text{邻域严格分开；}\qquad
\text{邻域严格分开}\Rightarrow\text{极限未必严格分开}.}
\]

还有一个配套反例用于防止反向误读：若
\[
f(x)=x,\qquad g(x)=0,\qquad x\to0,
\]
则
\[
\lim f=\lim g=0,
\]
但任意双侧去心邻域内都有 $x>0$ 与 $x<0$，所以既不能由“极限相等”推出最终 $f\ge g$，也不能推出最终 $f\le g$。\kw{极限相等时没有正间隙，邻域次序通常是不确定的。}
\end{corebox}

更一般地，\textbf{先假设相关极限存在}：
\[
f(x)\to A,\qquad g(x)\to B.
\]
此时大小关系才有两条方向完全不同的命题：
\begin{enumerate}
  \item \textbf{极限值有严格间隙 $\Rightarrow$ 邻域最终保持严格次序。} 若 $A>B$，令 $h=f-g$，则 $h\to A-B>0$，由保号性知充分靠近时 $f(x)>g(x)$。
  \item \textbf{邻域最终保持非严格次序 $+$ 两极限存在 $\Rightarrow$ 取极限后仍保持非严格次序。} 若充分靠近时 $f(x)\ge g(x)$，则必有 $A\ge B$。否则若 $A<B$，第一条应用于 $g-f$ 会推出充分靠近时 $g>f$，与题设矛盾。
\end{enumerate}

\begin{warnbox}{严格不等式经过极限通常会降级}
在已经确认两极限存在的前提下，若充分靠近时 $f(x)>g(x)$，仍只能推出 $A\ge B$，不能保证 $A>B$。例如 $x\to0$ 时
\[
f(x)=x^2>0=g(x),
\qquad
\lim f=\lim g=0.
\]
反过来，$A=B$ 也不能决定邻域中的大小方向；$f=-x^2$、$g=0$ 就给出相反次序。要把严格性传到极限，必须保留一个不会消失的统一间隙，例如存在 $\eta>0$ 使最终 $f\ge g+\eta$。
\end{warnbox}

\subsection{极限四则运算：先有部件极限，才能无损拼装}

设 $f(x)\to A$、$g(x)\to B$，其中 $A,B\in\mathbb R$。则有限极限对和、差、积封闭：
\[
\lim(f\pm g)=A\pm B,
\qquad
\lim(fg)=AB.
\]
当 $B\ne0$ 时，由极限保号/离零性可知 $g$ 在某个去心邻域内不为零，从而
\[
\lim\frac{f}{g}=\frac AB.
\]
这里“两个部件极限存在”是\kw{资格}，不是形式装饰。若一个部件极限不存在，不能机械把“$\lim$”拆进表达式。

反过来，四则关系有时能用于\kw{存在性反证}。若 $\lim f$ 存在而 $\lim g$ 不存在，则 $\lim(f+g)$ 不可能存在；否则由
\[
g=(f+g)-f
\]
会推出 $g$ 也有极限，矛盾。但乘积不能照抄这个逻辑：$g$ 不收敛而 $f\to0$ 时，$fg$ 仍可能因包络压制而收敛。

因此“存在 $\pm$ 不存在”在加减结构中可由可逆恒等式判断；“不存在 $\times$ 不存在”或“零 $\times$ 不存在”则必须回到具体尺度，不能靠标签运算。

四则运算还有两个常用的\kw{反向存在性推论}。若商本身在相应去心邻域有定义且
\[
\lim\frac{f}{g}=C\in\mathbb R,
\qquad
\lim g=0,
\]
则由 $f=(f/g)g$ 得
\[
\lim f=0.
\]
反过来，若
\[
\lim\frac{f}{g}=C\ne0,
\qquad
\lim f=0,
\]
则 $f/g$ 最终离零，从
\[
g=\frac{f}{f/g}
\]
可得
\[
\lim g=0.
\]
这两条都来自已经合法的四则拼装，不是新的极限口诀；若商极限不存在、商在邻域中无定义，或分母极限为零却准备直接套商法则，都必须退回资格检查。

\begin{corebox}{次序传递的三道门：存在性 $\to$ 方向 $\to$ 严格性}
\[
\boxed{\text{存在性与共同定义域}\to\text{传递方向}\to\text{严格性}}
\]

\textbf{第一道门：存在性与共同定义域。} 从邻域性质推极限结论时，必须先证明相关极限存在；从极限推邻域性质时，先取各自合法去心邻域的交集。

\textbf{第二道门：方向。} 从极限到邻域，需要极限值之间存在严格间隙；从邻域到极限，在极限已经存在时只能保留次序。

\textbf{第三道门：严格性。} 取极限时严格次序可能压到边界而退化成等号。于是“有没有统一正间隙”只在通过前两道门之后才是正确问题。
\end{corebox}

% ============================================================
\section{存在性：证明要控制全部路径，否定只需制造一次冲突}
% ============================================================

\subsection{直接控制与夹逼}

若能找到随 $t\to0$ 而趋于 $0$ 的控制函数 $\omega(t)$，使
\[
\lvert f(x)-L\rvert\le \omega(\lvert x-a\rvert),
\]
则只要把输入邻域缩到足以令 $\omega(\lvert x-a\rvert)<\varepsilon$，极限便成立。

夹逼定理是这一机制的有序版本：若
\[
g(x)\le f(x)\le h(x),\qquad g(x),h(x)\to L,
\]
则中间对象没有剩余自由度，只能趋于 $L$。

特别地，若 $u(x)$ 有界、$v(x)\to0$，则
\[
\lvert u(x)v(x)\rvert\le M\lvert v(x)\rvert\to0.
\]
这里真正起作用的不是“振荡自动消失”，而是\kw{振幅包络被压到零}。

\subsection{左右冲突与归结原则}

函数极限存在时，沿任何合法数列 $x_n\to a$（且 $x_n\ne a$）都必须得到同一函数值极限。反过来，这一性质对所有合法数列成立，也能推出函数极限存在。这就是海涅定理的接口。

因此否定极限通常比证明极限便宜：
\begin{itemize}
  \item 找到左右极限不等；
  \item 找到两条趋近数列，其函数值趋向不同常数；
  \item 找到一条趋近路径使函数值无界，另一条保持有限；
  \item 证明输出在两个分离区间间持续振荡。
\end{itemize}

\begin{corebox}{存在性证明与不存在证明不对称}
证明存在必须控制所有允许输入；检验有限条路径相同不能证明存在。否定存在只需一对合法路径冲突。这条不对称在多元极限中会更强：两条路径相同仍几乎没有证明力，两条路径不同却足以结束。
\end{corebox}

\subsection{贯穿母例：振荡信号与收缩包络}

定义
\[
F_p(x)=
\begin{cases}
\lvert x\rvert^p\sin\dfrac1x,&x\ne0,\\
c,&x=0.
\end{cases}
\]

当 $p=0$ 时，取
\[
x_n=\frac{1}{\frac\pi2+2n\pi},\qquad
y_n=\frac{1}{\frac{3\pi}{2}+2n\pi},
\]
两列都趋于 $0$，但 $F_0(x_n)=1$、$F_0(y_n)=-1$，所以极限不存在。

当 $p>0$ 时，
\[
\lvert F_p(x)\rvert\le \lvert x\rvert^p\to0,
\]
故由夹逼定理
\[
\lim_{x\to0}F_p(x)=0.
\]
此时函数在 $0$ 连续当且仅当 $c=0$。

这个母例同时暴露一个尺度陷阱：虽然 $F_p(x)=O(\lvert x\rvert^p)$，却没有
\[
F_p(x)\sim \lvert x\rvert^p,
\]
因为二者之比是 $\sin(1/x)$，没有极限。\kw{被某一阶控制}不等于\kw{拥有稳定的同阶主部}。

% ============================================================
\section{主导尺度：极限计算不是变形竞赛，而是信息精度管理}
% ============================================================

\subsection{先把趋近点归一化}

有限点 $x\to a$ 时令 $h=x-a$，把问题转成 $h\to0$；无穷远问题若适合倒代换，可令 $t=1/x$。归一化不是目的，它只是让所有局部尺度都用同一把尺比较。

\subsection{等价无穷小就是相对误差趋零}

若 $f(x),g(x)\to0$ 且去心邻域内 $g(x)\ne0$，则
\[
f(x)\sim g(x)
\iff \frac{f(x)}{g(x)}\to1
\iff f(x)=g(x)\bigl(1+o(1)\bigr).
\]

这不是“两个量相等”，而是说用 $g$ 替代 $f$ 时，\kw{相对误差}趋于零。更一般地：
\begin{align*}
f=o(g)&\iff \frac{f}{g}\to0,\\
f=O(g)&\iff \frac{f}{g}\text{ 在相关邻域有界}.
\end{align*}

信息强度依次为
\[
\text{恒等式} > \text{带余项展开} > \text{等价} > \text{同阶/大 }O > \text{仅有界}.
\]
使用较弱信息可以换取速度，但不能从较弱关系反推出较强关系。

\subsection{展开中心先于展开公式：真正趋零的是“相对中心的增量”}

局部公式从来不是因为题面写着 $x\to0$ 就自动可用。先问\kw{被展开的内部变量趋向哪里}，再把它改写成相对该中心的小增量。

例如 $x\to0$ 时
\[
e^{1+x+\frac12x^2}-1\to e-1\ne0,
\]
因此不能把整个式子套成 $e^u-1\sim u$。真正的小量是 $u=x+\frac12x^2\to0$，所以若需要研究围绕 $e$ 的变化，应写成
\[
e^{1+u}-e=e\bigl(e^u-1\bigr)\sim eu.
\]
同理，$\ln u\sim u-1$ 的资格是 $u\to1$；$e^u-1\sim u$ 的资格是 $u\to0$。\kw{展开中心错位会让后续阶数判断从第一步就失效。}

\subsection{为什么乘除可替换，而加减会失真？}

若 $f\sim g$，则 $f=g(1+o(1))$。在乘除结构中，这个相对误差通常仍是 $1+o(1)$，所以主部得以保留。

但若两个主部相减，低阶项可能完全抵消：
\[
\sin x\sim x,
\]
却不能据此把 $x-\sin x$ 写成 $x-x=0$。被消掉的正是当前已知的全部主部，答案由下一阶决定。此时必须切换到专题03的有限泰勒模型，或者使用能够保留下一阶信息的恒等变形。

\begin{methodbox}{尺度账本}
每次近似都记录三件事：
\begin{enumerate}
  \item 当前保留到哪一阶；
  \item 加减后这一阶是否抵消；
  \item 目标分母或比较尺度要求保留到哪一阶。
\end{enumerate}
若主部抵消，不能继续沿用原精度；必须升级信息等级。
\end{methodbox}

\subsection{未定式只报告“当前信息不足”}

$0/0$、$\infty/\infty$、$0\cdot\infty$、$\infty-\infty$、$1^\infty$、$0^0$、$\infty^0$ 不是答案类型，而是说当前粗尺度还无法决定结果。它们共同表达的是：\textbf{只知道各部件各自趋向哪里，还不足以决定组合后的极限；必须继续暴露相对速度、最低非零主部或更强约束。}

\begin{center}
\begin{tabularx}{\linewidth}{L{2.2cm}L{3.1cm}Y}
\toprule
\textbf{原始形态} & \textbf{先标准化成什么} & \textbf{为什么这样做}\\
\midrule
$0/0$ & 保持商型，先看能否约分、等价、局部展开或求导 & 两个量都趋零，关键是比较谁消失得更快\\
$\infty/\infty$ & 提最高增长尺度、倒数化或保持商型 & 关键是比较两者增长速度，而不是把“无穷”当数相除\\
$0\cdot\infty$ & $u/(1/v)$ 或 $v/(1/u)$ & 把乘积竞争改写为可比较的商型竞争\\
$\infty-\infty$ & 通分、有理化、提最大项或同尺度归一 & 差的结果由两项主部是否抵消决定\\
$1^\infty$ & 取对数：$v\ln u$ & 底数的一阶偏移与指数发散速度共同决定结果\\
$0^0,\ \infty^0$ & 取对数：$v\ln u$ & 幂型不确定性被翻译为乘积尺度问题\\
\bottomrule
\end{tabularx}
\end{center}

因此“识别未定式”不是做题终点，而是一道\kw{信息需求诊断}：当前还缺的是相对尺度、下一阶主部、统一上下界，还是导数/累积提供的另一种表示？

更精确地说：若底数 $u(x)>0$ 且 $u(x)\to0$，指数 $v(x)\to+\infty$，则 $u(x)^{v(x)}\to0$，所以这种合法的 $0^{+\infty}$ 情形不是未定式。若底数符号、实数幂定义域或指数方向不稳定，必须先回到专题01检查对象是否合法。

\subsection{标准化的四种基本变换：每一步都要知道保留了什么}

\begin{enumerate}
  \item \kw{恒等变形}：约分、通分、拆项、提公因子、三角恒等式。若在合法定义域内等价，它几乎不丢信息，优先级最高。
  \item \kw{有理化}：把根式差
  \[
  \sqrt{A(x)}-\sqrt{B(x)}
  \]
  改写为
  \[
  \frac{A(x)-B(x)}{\sqrt{A(x)}+\sqrt{B(x)}}.
  \]
  它的价值不是“看到根式就乘共轭”，而是把隐藏在平方根里的\textbf{差尺度}暴露出来。趋于负无穷时还必须单独处理 $\sqrt{x^2}=\lvert x\rvert$ 的符号。
  \item \kw{变量归一}：有限点令 $h=x-a$，无穷远可令 $t=1/x$ 或先提最高增长尺度。它把不同题面统一到 $h\to0$ 或一个明确尾部上，便于比较阶数。
  \item \kw{取对数}：只在幂型且底数保持为正时使用。若 $y=u(x)^{v(x)}$，则
  \[
  \ln y=v(x)\ln u(x).
  \]
  先求乘积极限，再用指数连续性恢复原极限。取对数改变的是表示，不允许忽略原表达式的正值与定义域资格。
\end{enumerate}

\subsection{离散最大尺度：有限个指数竞争时，$n$ 次根只留下最大者}

设 $c_1,\dots,c_m\ge0$ 为固定常数，记
\[
M=\max_{1\le j\le m}c_j.
\]
则
\[
M^n\le \sum_{j=1}^m c_j^n\le mM^n,
\]
因此
\[
M\le \sqrt[n]{\sum_{j=1}^m c_j^n}\le m^{1/n}M\to M.
\]
故
\[
\boxed{\lim_{n\to\infty}\sqrt[n]{c_1^n+\cdots+c_m^n}=\max_j c_j.}
\]

这不是新的孤立公式，而是主导尺度在离散指数结构中的极端版本：比最大底数小的项经过 $n$ 次幂后按指数级被压掉，开 $n$ 次根又把共同的指数尺度还原。若 $m=m_n$ 随 $n$ 增长，上界会多出 $m_n^{1/n}$；只有当
\[
m_n^{1/n}\to1
\]
时，原来的“最大项原则”仍可直接保留。若允许负底数，还必须先审查奇偶 $n$ 下实数定义与符号抵消，不能把非负情形无条件推广。

\begin{corebox}{母例：参数 $x$ 改变时，最大尺度会发生分区切换}
对固定 $x\ge0$，考虑
\[
A_n(x)=\sqrt[n]{1+x^n+\left(\frac{x^2}{2}\right)^n}.
\]
直接由最大尺度原则，
\[
\lim_{n\to\infty}A_n(x)=\max\left\{1,x,\frac{x^2}{2}\right\}.
\]
比较三个候选尺度便得到
\[
\lim A_n(x)=
\begin{cases}
1,&0\le x\le1,\\
x,&1\le x\le2,\\
\dfrac{x^2}{2},&x\ge2.
\end{cases}
\]
这个母例训练的不是“记三个区间”，而是：\textbf{先列候选主尺度，再找参数在哪里让主导者换位。}
\end{corebox}

\subsection{连乘开 $n$ 次根：真正被平均的是对数，不是因子本身}

若 $u_{n,k}>0$，则
\[
\left(\prod_{k=1}^{m_n}u_{n,k}\right)^{1/n}
=
\exp\left(\frac1n\sum_{k=1}^{m_n}\ln u_{n,k}\right).
\]
因此这类题的核心对象是\kw{平均对数尺度}。若右端和式可化为黎曼和，就转专题05处理累积；若只需比较是否趋于 $1$，则只需判断
\[
\frac1n\sum_{k=1}^{m_n}\ln u_{n,k}\to0.
\]

一个重要边界是：\textbf{整体乘积趋于 $0$，并不妨碍它的 $n$ 次根趋于 $1$。} 例如
\[
P_n=\prod_{k=1}^n\frac{2k-1}{2k}
=\frac{\binom{2n}{n}}{4^n}\to0,
\]
但由
\[
\frac1{2n+1}\le P_n\le1
\]
可得
\[
(2n+1)^{-1/n}\le P_n^{1/n}\le1,
\]
从而
\[
P_n^{1/n}\to1.
\]
真正决定 $n$ 次根的是 $\ln P_n$ 相对于 $n$ 的尺度，而不是 $P_n$ 本身是否趋零。

\subsection{中值定理的差值运输接口：外层相同、输入靠近时先搬运差}

若 $u(x),v(x)\to a$，且 $\Phi$ 在两者之间满足 Lagrange 中值定理资格，则对每个充分靠近且 $u(x)\ne v(x)$ 的 $x$，存在位于 $u(x)$ 与 $v(x)$ 之间的 $\xi_x$ 使
\[
\Phi(u(x))-\Phi(v(x))=\Phi'(\xi_x)\bigl(u(x)-v(x)\bigr).
\]
于是 $\xi_x\to a$；若进一步有 $\Phi'(t)\to\Phi'(a)$，便得到
\[
\frac{\Phi(u(x))-\Phi(v(x))}{u(x)-v(x)}\to\Phi'(a).
\]
这个接口的价值是把两个外层函数值的\kw{差}一次性运输成“局部变化率 $\times$ 输入差”，避免分别展开后再处理相消。中值定理本体与资格由专题04拥有；本册只使用它比较局部尺度。

\subsection{方法选择的知识层判据：不是题型表，而是“需要什么信息”}

极限计算的核心不是记“某类题用某招”，而是判断\textbf{当前表示缺哪一级信息}：

\begin{tabularx}{\linewidth}{L{3.0cm}L{3.0cm}Y}
\toprule
\textbf{当前缺失的信息} & \textbf{优先接口} & \textbf{典型原因}\\
\midrule
统一绝对误差上界 & 夹逼/有界量乘无穷小 & 振荡本身不可控，但振幅可被压到零\\
稳定最低阶主部 & 等价无穷小 & 乘除结构只需相对主部，不需要更高阶细节\\
主部抵消后的下一阶 & 专题03的有限泰勒 & 加减把现有主部消掉，低精度信息已经失效\\
商的相对变化率 & 专题04的洛必达 & 合法的 $0/0$ 或 $\infty/\infty$，且求导确实降低复杂度\\
固定基点的一阶变化率 & 专题03的导数定义 & 差商已围绕同一基点，可直接解释成局部一阶系数\\
离散和的连续累积极限 & 专题05的黎曼和接口 & 归一化后出现 $\frac1n f(k/n)$ 的网格累积结构\\
收缩积分区间的局部面积 & 专题05的变限积分接口 & 区间宽度与被积函数局部高度共同决定主部\\
\bottomrule
\end{tabularx}

\begin{corebox}{方法不会替你判断信息精度}
等价无穷小、洛必达、泰勒都可能算出同一道题，但它们保留的信息量不同。若一阶主部已经足够，泰勒是过度展开；若一阶主部相消，等价就是信息不足；若求导后结构更坏，洛必达即使形式可用也不是好的表示。真正的停止条件是：\textbf{目标极限已由当前信息唯一决定，而且剩余误差严格低于目标尺度。}
\end{corebox}

\subsection{几类特殊结构怎样进入正确归属专题}

旧题型笔记中的很多“极限方法”其实是跨专题接口，不能全部复制进专题02，但必须在这里保留路由逻辑：
\begin{itemize}
  \item \kw{差函数与参数配阶}：先在本册判断主部是否抵消；需要具体高阶系数时转专题03的泰勒模型。
  \item \kw{导数定义型极限}：本册负责识别“这是局部尺度问题”；固定基点差商的机制与左右导数边界归专题03。
  \item \kw{变上限积分极限}：本册负责比较收缩速度和主导尺度；“高度 $\times$ 宽度”、莱布尼茨求导与黎曼累积本体归专题05/H-B03。
  \item \kw{和式极限}：本册先判断是否存在统一尺度与误差压制；标准黎曼和及网格累积机制归专题05。
  \item \kw{增长阶比较}：$\ln^k x\ll x^\alpha\ll b^x$ 属于尺度锚点；若有差分抵消、振荡或符号分支，不能仅凭增长阶结束。
\end{itemize}

\subsection{离散差分接口：斯托尔茨定理}

数列比值有时不是靠连续化，而是靠“\textbf{累计量的比值由相邻增量比控制}”。设 $b_n$ 严格递增且 $b_n\to+\infty$。若
\[
\lim_{n\to\infty}\frac{a_{n+1}-a_n}{b_{n+1}-b_n}=L
\]
存在（允许扩展实数意义下的无穷极限），则斯托尔茨定理给出
\[
\lim_{n\to\infty}\frac{a_n}{b_n}=L.
\]

它的认知作用不是“把数列极限机械改成作差”，而是把 $b_n$ 看成不断累积的尺度：每新增一单位尺度时，$a_n$ 的局部增量比若稳定，长期累计的平均比值也被锁定。特别地，当 $b_n=n$ 时，问题从 $a_n/n$ 转成相邻差 $a_{n+1}-a_n$。

\begin{warnbox}{斯托尔茨不是离散版洛必达口诀}
必须先验证分母序列严格递增并趋于 $+\infty$；差分比只是满足这些资格后的接口。分母振荡、有界或差分比本身没有稳定极限时，不能因为“分子分母都趋于无穷”就强行作差。
\end{warnbox}

\subsection{工具的知识接口}

\begin{tabularx}{\linewidth}{L{2.5cm}YY}
\toprule
\textbf{工具} & \textbf{它提供的信息} & \textbf{调用边界}\\
\midrule
恒等变形/有理化 & 不损失信息地改变表示 & 先维护定义域与非零条件\\
夹逼/有界量乘无穷小 & 给出统一绝对误差包络 & 需要真实上下界或有界性\\
等价无穷小 & 提取稳定相对主部 & 加减相消时信息不足\\
洛必达 & 把合法商型的极限转移到导数比，用局部变化率比较原商 & 只对满足条件的 $0/0$ 或 $\infty/\infty$；机制与资格链归专题04，每次连续使用都必须重新验型与验资格；导数比极限不存在\textbf{不能反推}原商极限不存在\\
有限泰勒 & 给出指定阶主部与余项 & 局部模型归专题03；必须按目标阶数截断\\
黎曼和/变限积分 & 把离散和或收缩区间翻译成累积 & 累积对象归专题05\\
斯托尔茨 & 用相邻增量比控制累计量比值 & 分母序列须严格递增且无界，差分比极限须存在\\
\bottomrule
\end{tabularx}

% ============================================================
\section{连续：把邻域趋势与点值接成同一个对象}
% ============================================================

\subsection{连续的三种等价语言}

函数 $f$ 在 $a$ 连续，等价于
\[
\lim_{x\to a}f(x)=f(a).
\]
它包含三个资格：$f(a)$ 有定义、双侧极限存在、极限等于点值。

用增量语言，令 $\Delta x=x-a$，则
\[
\Delta x\to0\quad\Longrightarrow\quad
\Delta y=f(a+\Delta x)-f(a)\to0.
\]

用 $\varepsilon$--$\delta$ 语言，连续只是把极限定义中的目标 $L$ 指定成 $f(a)$，并允许 $x=a$：
\[
\forall\varepsilon>0,\ \exists\delta>0,
\quad \lvert x-a\rvert<\delta\Longrightarrow \lvert f(x)-f(a)\rvert<\varepsilon.
\]

\subsection{单侧连续与区间连续：连续性必须服从合法输入集合}

右连续与左连续分别是
\[
\lim_{x\to a^+}f(x)=f(a),
\qquad
\lim_{x\to a^-}f(x)=f(a).
\]
若 $a$ 是定义域内部点，则
\[
f\text{ 在 }a\text{ 连续}
\iff
f\text{ 在 }a\text{ 左连续且右连续}.
\]
这里的“两侧”不是形式要求，而是当前定义域确实允许的两条趋近通道。

因此闭区间上的连续应写成：$f$ 在 $(a,b)$ 每一点连续，并且在左端点 $a$ 右连续、在右端点 $b$ 左连续，记为
\[
f\in C[a,b].
\]
这与前文的合法输入集合完全一致：端点没有定义域外侧的输入，不需要虚构另一侧极限。闭区间连续以后能产生哪些全局性质，由专题04负责。

\subsection{连续函数保号只是“极限保号 + 点值接缝”}

若 $f$ 在 $a$ 连续且 $f(a)>0$，则
\[
\lim_{x\to a}f(x)=f(a)>0.
\]
由前面的极限保号性，存在 $\delta>0$，使充分靠近 $a$ 的输入都有 $f(x)>0$；而 $x=a$ 本身又满足 $f(a)>0$，故整个小邻域内保持正号。$f(a)<0$ 同理。

所以连续函数保号不需要再记成一条孤立定理：它只是
\[
\boxed{\text{点值有严格符号}\xrightarrow{\text{连续接缝}}\text{极限有同一严格符号}\xrightarrow{\text{极限保号}}\text{邻域最终同号}}.
\]
如果只知道 $f(a)=0$，这条链立即失效；连续函数可以在该点穿过零点、接触零点，也可以在两侧同号。

\subsection{运算封闭性是带资格的}

若 $f,g$ 在 $a$ 连续，则 $f\pm g$、$fg$ 连续；只有当 $g(a)\ne0$ 时，$f/g$ 才连续。若 $g$ 在 $a$ 连续且 $\varphi$ 在 $g(a)$ 连续，则 $\varphi\circ g$ 在 $a$ 连续。

复合的正确观察方式不是只看“内外层是否各自有坏点”，而是追踪：
\[
x\to a\quad\Longrightarrow\quad g(x)\to b
\quad\Longrightarrow\quad \varphi(g(x))\to\varphi(b).
\]
内层的实际输出路径可能避开外层坏点；非单射外层也可能把不同的左右极限压成同一个值。

\subsection{间断分类就是定位接缝在哪一层失败}

记
\[
L^- = \lim_{x\to a^-}f(x),\qquad
L^+ = \lim_{x\to a^+}f(x).
\]

\begin{tabularx}{\linewidth}{L{2.2cm}L{4.1cm}YY}
\toprule
\textbf{类型} & \textbf{判据} & \textbf{接缝故障} & \textbf{单点修复能力}\\
\midrule
可去间断 & $L^-=L^+=L$，但 $f(a)$ 缺失或不等于 $L$ & 邻域已一致，点值接错 & 令 $f(a)=L$ 即可\\
跳跃间断 & $L^-,L^+$ 有限但不相等 & 两侧稳定到不同状态 & 改一个点无效\\
无穷间断 & 至少一侧趋于 $\pm\infty$ & 输出逃离所有有限范围 & 改一个点无效\\
振荡间断 & 至少一侧持续振荡而无极限 & 邻域内没有唯一输出状态 & 改一个点无效\\
\bottomrule
\end{tabularx}

“第一类间断”要求左右极限都是有限值；至少一侧没有有限极限则进入第二类。定义域端点只按合法一侧讨论，不机械归入双侧间断分类。

\subsection{修复、抵消与信息丢失}

连续性的封闭定理是\kw{正向充分条件}，不能无条件反推：
\begin{itemize}
  \item 两个不连续函数可能相加后抵消成连续函数；
  \item 趋于零的因子可以压制有界振荡；
  \item 平方或绝对值会丢失符号，可能把左右极限 $A\ne B$ 压成同一输出；
  \item 非零连续因子不会修复另一个函数的间断，因为局部可以除回去。
\end{itemize}

例如若 $f$ 在 $a$ 有有限单侧极限 $A\ne B$，外层 $\varphi$ 在 $A,B$ 连续，则
\[
\lim_{x\to a^-}\varphi(f(x))=\varphi(A),\qquad
\lim_{x\to a^+}\varphi(f(x))=\varphi(B).
\]
只有当外层把两者映成同一值，并且点值也接上时，跳跃才可能被隐藏。这里发生的是\kw{输出信息压缩}，不是原函数突然连续。

\subsection{扩展：连续函数列的点态极限为什么可能断掉？}

若每个 $f_n$ 都连续，而只知道对每个固定 $x$ 都有
\[
f_n(x)\to f(x),
\]
这只是\kw{点态收敛}：允许达到给定误差所需的 $N$ 随 $x$ 改变。不同空间点可以以完全不同的速度收敛，因此无法自动把“每个 $f_n$ 的连续控制”统一传给极限函数。

例如在 $[0,1]$ 上令 $f_n(x)=x^n$。每个 $f_n$ 都连续，但点态极限为
\[
f(x)=\begin{cases}
0,&0\le x<1,\\
1,&x=1,
\end{cases}
\]
在 $x=1$ 不连续。

若进一步有\kw{一致收敛}，即对每个 $\varepsilon>0$ 都能找到一个与 $x$ 无关的统一 $N$，使全体 $x$ 同时满足 $\lvert f_n(x)-f(x)\rvert<\varepsilon$，那么连续性才会传给极限函数。因此在连续函数列的语境中：\textbf{若点态极限已经不连续，则这列函数不可能在该集合上一致收敛。}

这条结论在本册只作为扩展，用来解释“连续对象逐点取极限为什么会产生间断”；一般函数列的一致收敛理论不扩入数学一核心主干。

% ============================================================
\section{递推数列：把极限放回一条离散轨道}
% ============================================================

对一阶自治递推
\[
x_{n+1}=\varphi(x_n),
\]
方程 $L=\varphi(L)$ 只给出\kw{可能的终点}。正确生命周期是
\begin{center}
\begin{tikzpicture}[
  node distance=0.8cm and 0.8cm,
  box/.style={draw=deep,rounded corners=1.8mm,minimum width=3.8cm,minimum height=1.05cm,align=center,fill=soft,font=\small,inner sep=4pt},
  succbox/.style={draw=greenish,rounded corners=1.8mm,minimum width=3.8cm,minimum height=1.05cm,align=center,fill=white,font=\small,inner sep=4pt,thick},
  arr/.style={-{Latex[length=2.2mm,width=1.5mm]},thick,draw=deep},
  warnarr/.style={-{Latex[length=2.2mm,width=1.5mm]},thick,draw=warn,densely dashed}
]
% 第一排：轨道准备与收敛保证
\node[box] (init) at (0,0)    {\textbf{1. 合法初值}\\[2pt]\scriptsize 确定初值点 $x_1\in I$};
\node[box] (inv)  at (4.8,0)  {\textbf{2. 不变区间}\\[2pt]\scriptsize 验证映射保持 $\varphi(I)\subseteq I$};
\node[box] (eng)  at (9.6,0)  {\textbf{3. 收敛引擎}\\[2pt]\scriptsize 单调有界 / 压缩映射};

% 第二排：极限确立与不动点筛选
\node[succbox] (lim)  at (9.6,-2.3) {\textbf{\color{greenish}4. 极限确立}\\[2pt]\scriptsize 证明存在 $\lim x_n = L$};
\node[box]     (cont) at (4.8,-2.3) {\textbf{5. 连续传递}\\[2pt]\scriptsize 传入方程 $L = \varphi(L)$};
\node[box]     (filt) at (0,-2.3)   {\textbf{6. 候选根筛选}\\[2pt]\scriptsize 结合轨道范围排除伪根};

\draw[arr] (init) -- (inv);
\draw[arr] (inv) -- (eng);
\draw[arr] (eng) -- (lim);
\draw[arr] (lim) -- (cont);
\draw[arr] (cont) -- (filt);

% 警示旁路：未证收敛直接设极限
\draw[warnarr] (inv.south) -- node[fill=white,inner sep=2.5pt,font=\scriptsize\color{warn},align=center] {\textbf{致命陷阱}：未证收敛直接设极限求解\\（必要非充分：无法排除发散、二周期或多根错选）} (cont.north);

\end{tikzpicture}
\end{center}

\subsection{先保证轨道存在并留在可控区域}

找区间 $I$，验证
\[
x_1\in I,\qquad \varphi(I)\subseteq I.
\]
这一步同时保证根式、分母等始终合法，并为有界性、导数界或符号分析提供统一区域。局部条件 $\lvert\varphi'(L)\rvert<1$ 只描述足够靠近 $L$ 后的吸引性，不能代替“初始轨道最终进入吸引域”的证明。

\subsection{轨道单调看迭代增量，不是只看映射单调}

因为
\[
x_{n+1}-x_n=\varphi(x_n)-x_n,
\]
真正决定轨道增减的是 $h(x)=\varphi(x)-x$ 在轨道所在区间的符号。$\varphi$ 递增只说明次序在一次迭代后被保持，并不自动说明 $x_{n+1}\ge x_n$。

若 $\varphi$ 递减，轨道常在不动点两侧交替。此时研究 $\varphi\circ\varphi$ 与奇偶子列；即使两条子列分别收敛，也必须排除
\[
L_2=\varphi(L_1),\qquad L_1=\varphi(L_2),\qquad L_1\ne L_2
\]
这种二周期。反例 $x_{n+1}=1-x_n$ 说明“有唯一不动点”也不能自动排除二周期轨道。

\subsection{离散母例：牛顿递推逼近根号二}

设 $x_1>0$，
\[
x_{n+1}=\frac12\left(x_n+\frac2{x_n}\right).
\]

首先由均值不等式
\[
x_{n+1}\ge\sqrt{x_n\cdot\frac2{x_n}}=\sqrt2,
\]
故从 $x_2$ 起轨道始终合法且有下界。又当 $x_n\ge\sqrt2$ 时，
\[
x_{n+1}-x_n=\frac{2-x_n^2}{2x_n}\le0,
\]
所以 $\{x_n\}_{n\ge2}$ 单调递减且有下界，极限存在。

设极限为 $L>0$。在确认极限存在且递推映射连续后，才可把极限传入递推式：
\[
L=\frac12\left(L+\frac2L\right)
\quad\Longrightarrow\quad L^2=2
\quad\Longrightarrow\quad L=\sqrt2.
\]

误差还有精确恒等式
\[
x_{n+1}-\sqrt2
=\frac{{(x_n-\sqrt2)}^2}{2x_n},
\]
说明靠近终点后误差近似平方缩小。这个例子把“候选终点”和“轨道确实到达终点”严格分开。

\begin{warnbox}{先设极限再解不动点方程，只完成了必要条件}
若未证明收敛，写 $L=\varphi(L)$ 只能说“任何可能的有限极限必须是不动点”。它不能排除发散、振荡、二周期，也不能决定多个不动点中轨道会到哪一个。
\end{warnbox}

\subsection{双数列均值递推：先建立夹链，再让两个轨道合流}

有一类递推不是单个 $x_{n+1}=\varphi(x_n)$，而是两个量互相取均值，例如算术--几何平均型结构。此时最有用的对象不是分别猜两个极限，而是寻找一条保持的\kw{夹链}：
\[
b_n\le b_{n+1}\le a_{n+1}\le a_n.
\]
一旦这条链成立，便立即得到：$\{a_n\}$ 单调递减且有下界，$\{b_n\}$ 单调递增且有上界，所以分别存在极限
\[
a_n\to A,\qquad b_n\to B,
\]
并且始终 $B\le A$。

最后才把递推关系传入极限。若递推的连续性进一步强迫 $A=B$，两条轨道便\kw{合流到同一个终点}。这比“先设两个极限再解方程”多了一步真正的存在性证明。

典型地，若
\[
a_{n+1}=\frac{a_n+b_n}{2},\qquad
b_{n+1}=\sqrt{a_nb_n},\qquad 0<b_1\le a_1,
\]
由算术--几何平均不等式
\[
b_n\le b_{n+1}\le a_{n+1}\le a_n,
\]
两列分别收敛；再令 $n\to\infty$ 得
\[
A=\frac{A+B}{2},
\]
故 $A=B$。

\begin{corebox}{双轨递推的母模型}
\[
\boxed{\text{排序/不变量}\to\text{夹链}\to\text{一增一减}\to\text{各自收敛}\to\text{极限方程迫使合流}}
\]
看到“两个数互相取算术/几何/调和均值”时，先找次序是否被保持以及两条轨道是否互相夹住，而不是先做复杂通项展开。
\end{corebox}

\subsection{扩展：压缩与误差传播}

若 $I$ 是完备的闭区间、$\varphi(I)\subseteq I$，且存在 $0\le q<1$ 使
\[
\lvert\varphi(x)-\varphi(y)\rvert\le q\lvert x-y\rvert\qquad(x,y\in I),
\]
则误差满足几何压缩，轨道收敛到 $I$ 内唯一不动点。这一压缩映射语言可作为递推题的强工具，但复杂动力系统与完整巴拿赫压缩原理不进入数学一主干。

% ============================================================
\section{概念边界：真正判据、题目信号与错误后果}
% ============================================================

\noindent
\begin{tabularx}{\linewidth}{L{2.1cm}>{\centering\arraybackslash}p{0.45cm}L{2.1cm}L{5.8cm}Y}
\toprule
\textbf{概念 A} & & \textbf{概念 B} & \textbf{为什么容易混 / 真正判据与题目信号} & \textbf{混淆后果}\\
\midrule
极限 & $\neq$ & 函数值 & 极限看去心邻域，函数值看点上定义；补定义、分段点和可去间断是信号。 & 直接代入，把一个点的信息当成邻域信息\\
有限极限 & $\neq$ & 无穷大记号 & 有限极限稳定到实数；$+\infty$ 表示最终超过任意高度，是扩展发散状态。 & 把发散方向当作普通数参与四则运算\\
有界 & $\neq$ & 收敛 & 有界只把输出关在某个范围；收敛还要求最终进入任意小的同一邻域。振荡是典型信号。 & 由有界直接宣布极限存在\\
若干子列同极限 & $\neq$ & 原列收敛 & 原列收敛会让任意子列继承同一极限；反向只有当这些子列在充分靠后时覆盖全部原下标才成立。奇偶子列覆盖全部下标，$3n$ 与 $3n+1$ 则漏掉 $3n+2$。 & 看见两条“漂亮子列”同极限就忽略未覆盖的尾部\\
$a_n-b_n\to0$ & $\neq$ & 两列都收敛 & 差趋零只说明两列渐近合流；若其中一列已有极限，另一列继承同一极限，但两列也可能同步振荡。 & 把“彼此接近”误当成“共同到达某个终点”\\
邻域次序 & $\neq$ & 极限次序 & 邻域内 $f\ge g$ 或 $f>g$ 只描述最终点态关系；要推出 $\lim f\ge\lim g$，必须先保证两个极限存在。反向若两极限存在且有严格间隙，先取共同去心邻域再保序。 & 跳过存在性/共同定义域，直接在函数值与极限之间搬运不等号\\
左右极限 & $\neq$ & 双侧极限 & 双侧存在要求两侧都存在且相等；分段、绝对值、定义域边界必须拆侧。 & 漏掉单侧冲突或要求不存在的另一侧\\
$O(g)$ & $\neq$ & $\sim g$ & $O(g)$ 只要求比值有界；等价要求比值趋于 $1$。有界振荡乘包络是典型反例。 & 把尺度上界误当稳定主部\\
等价 & $\neq$ & 相等 & 等价只控制趋近过程中的相对误差；加减相消会暴露被丢掉的高阶项。 & 在相消结构中直接替换，得到错误零值\\
极限候选 & $\neq$ & 极限存在 & 代入、数值观察或不动点方程只能产生候选；存在性还要控制全部路径/尾部。 & 先设极限后忘记证明收敛\\
映射单调 & $\neq$ & 轨道单调 & 递推轨道增减看 $\varphi(x_n)-x_n$；映射单调只说明次序如何传播。 & 用 $\varphi'\ge0$ 直接证明数列递增\\
不动点唯一 & $\neq$ & 轨道收敛 & 唯一不动点不排除发散或二周期；还需不变区间与收敛发动机。 & 解出唯一根就结束递推证明\\
连续 & $\neq$ & 可导 & 连续只把邻域趋势接到点值；可导还要求统一的一阶局部主部，由专题03拥有。 & 从连续直接推出导数存在\\
点态极限 & $\neq$ & 一致极限 & 点态控制阈值可依赖空间点；一致控制要求一个阈值对全域有效。 & 把连续函数列的点态极限无条件判为连续\\
\bottomrule
\end{tabularx}

% ============================================================
\section{最小调用接口：把学生送进正确模型即可}
% ============================================================

本册规范正文的责任到“识别当前在处理哪一种极限机制、知道缺什么信息、知道何时转交”即停止。完整题型树、局部技巧和考场状态机由同目录《极限计算》《数列极限的结构转译》《数列收敛性与逻辑判定》《连续与间断》《递推数列》承担。

\subsection{题目信号}

看到下列对象时，应回到本册：
\begin{itemize}
  \item $x\to a^\pm$、$x\to a$、$x\to\pm\infty$ 或 $n\to\infty$；
  \item 奇偶子列、模类子列、抽取子列或 $a_n-b_n\to0$ 等需要判断“局部观察能否控制全部尾部”的结构；
  \item 分段、绝对值、取整、振荡、定义域端点等可能改变合法输入集合或路径一致性的结构；
  \item $0/0$、$\infty/\infty$ 等未定式，或根式差、指数差、幂型等把相对尺度隐藏起来的表示；
  \item 连续/间断问题中的左右极限与点值接缝；
  \item 递推数列中“候选不动点”与“轨道是否真的收敛”的区分。
\end{itemize}

\subsection{第一观察层}

进入本册后只先确认三件事：
\begin{enumerate}
  \item \textbf{允许怎样趋近？} 定义域是否迫使拆左右、单侧或子列？若只观察若干子列，它们是否覆盖原列全部尾部？
  \item \textbf{极限是否已经有存在性资格？} 是否能立即找到路径冲突、振荡或无界变号？“两列之差趋零”这类关系是在传递已有存在性，还是被误用来制造存在性？
  \item \textbf{当前信息精度够不够？} 最低阶主部是否已经唯一决定答案，还是加减相消后必须升级到下一阶？
\end{enumerate}

这三问给出机制入口，但不规定一条所有题都必须执行到底的算法。

\subsection{退出与转交边界}

\begin{itemize}
  \item 只需相对主部时，可在本册调用等价、夹逼等尺度机制；主部相消并需要具体高阶系数时转专题03的有限泰勒；
  \item 差值适合用中值点运输，或商型适合洛必达时，转专题04检查相应定理资格；
  \item 出现黎曼和、收缩积分区间或变上限积分时，转专题05；
  \item 固定基点差商已经成为导数定义时，转专题03；
  \item 连续接缝已经确认后，不在本册继续越级推可导；递推极限已证存在后，才允许用不动点方程筛选终点。
\end{itemize}

% ============================================================
\section{认知边界与失败路径}
% ============================================================

\begin{tabularx}{\linewidth}{L{3.2cm}YY}
\toprule
\textbf{错误直觉} & \textbf{为什么失效} & \textbf{应回到的机制}\\
\midrule
分段点只代入某一侧 & 双侧极限要求所有合法路径一致，点值又属于更后的接缝层 & 合法输入 / 路径一致\\
有界就等于收敛 & 有界只限制范围，不能提供唯一稳定状态 & 夹逼、归结原则或冲突子列\\
邻域内 $f>g$ 就直接搬到极限 & 邻域次序本身不制造两个极限的存在性，严格号取极限还可能退化 & 极限存在性 + 次序传递边界\\
等价替换后把差式变成 $0$ & 当前最低阶信息已被相消，答案由被丢掉的下一阶决定 & 尺度账本；必要时转专题03\\
连续使用洛必达 & 每次得到的新商都是新的极限问题，且定理是单向充分路径 & 专题04的资格链\\
只看 $f(a)$ 判断连续 & 点值不能代表去心邻域行为 & 左右极限 + 点值接缝\\
复合函数只看外层是否连续 & 内层实际输出路径可能触及不同分支，外层也可能压掉信息 & 复合极限的路径追踪\\
递推题先解 $L=\varphi(L)$ & 不动点方程只筛候选终点，不证明轨道到达 & 不变区间 + 收敛发动机\\
奇偶子列各自收敛就结束 & 两个子列极限可能构成二周期而不相等 & 二步映射与二周期排除\\
\bottomrule
\end{tabularx}

% ============================================================
\section{考纲覆盖与归属边界}
% ============================================================

按考纲查漏，本册覆盖：数列与函数极限、数列极限唯一性与子列继承/覆盖反推、左右极限、极限性质与运算、无穷小/无穷大及阶比较、七类未定式判别与标准化、斯托尔茨差分比接口、夹逼与单调有界准则、渐近合流（差趋零）的存在性边界、连续与间断、初等函数连续性。

具体计算工具不在本册重复拥有：有限泰勒与导数定义归专题03，洛必达及其中值定理资格归专题04，黎曼和与变限积分归专题05，级数尾部归专题10。本册只解释“当前缺哪一级信息、为什么应转交”；具体题面触发、第一动作与退出条件由同目录训练 Markdown 维护。来源材料的逐项吸收关系由证据层的来源差异记录维护，不在规范正文维护第二份迁移台账。

% ============================================================
\section{一页压缩：忘掉公式后怎样重建？}
% ============================================================

\begin{corebox}{一句话}
极限是在指定趋近方式下，用一个阈值统一控制所有足够靠近的合法输入；主导尺度记录控制所需的信息精度；连续只是再把这个邻域极限接到点值。
\end{corebox}

\subsection{母图}
\[
\boxed{
\text{趋近规范}
\to \text{合法邻域/尾部}
\to \text{统一误差包络}
\to \text{所有路径一致}
\to \text{与点值接缝}}
\]

\subsection{四个生成原则}
\begin{enumerate}
  \item \kw{存在性原则}：证明要控制全部允许输入；否定只需制造一对合法冲突。任何结论一旦出现极限，先问“这个极限是谁保证存在的？”
  \item \kw{次序原则}：邻域与极限之间传递不等式时，先过存在性/共同定义域门；极限的严格间隙可以保留到邻域，而邻域的严格次序取极限时可能降级为非严格次序。
  \item \kw{尺度原则}：近似是在丢信息；主部抵消时必须升级精度，不能继续使用已失效的等价。
  \item \kw{接缝原则}：极限、点值、连续是三层信息；改变单点值只能修复可去间断，不能修复邻域结构。
\end{enumerate}

\subsection{重建问题}
\begin{enumerate}
  \item 当前允许怎样趋近？定义域是否迫使我拆左右或只看单侧？若在看子列，它们只是抽样观察，还是已经覆盖全部尾部？
  \item 结论中出现的每个极限，题设或前面的推理是否已经保证它存在？比较两个函数时，是否已经取得共同合法去心邻域？两列之差趋零时，是谁先提供了存在性？
  \item 我是在产生候选值，还是已经证明所有路径一致？若若干子列同极限，还有没有未被控制的下标类？
  \item 当前近似保留到哪一阶？加减后主部是否抵消？
  \item 外层运算保留了哪些信息，又压掉了哪些差异？
  \item 连续题是否核对了左右极限与点值三者？
  \item 递推题是否先证明轨道存在并收敛，再使用不动点方程？
  \item 当前机制属于专题02，还是应该切换到专题03/04/05/10？
\end{enumerate}

\begin{center}\rule{0.5\linewidth}{0.5pt}\end{center}

\appendix
\section{尺度锚点与考纲速查}

本附录只保存可由母模型调用的尺度锚点，不替代正文中的存在性、定义域与相消检查。使用前先把局部变量归一化为 $h=x-a\to0$。

\subsection{两个重要极限}

\[
\lim_{h\to0}\frac{\sin h}{h}=1,
\qquad
\lim_{h\to0}{(1+h)}^{1/h}=e\quad(1+h>0).
\]

第二个极限也可写成
\[
\lim_{u\to\infty}{\left(1+\frac1u\right)}^u=e.
\]
它们不是孤立口诀：前者给三角函数的局部线性尺度，后者给“单位底数的微小偏移 $\times$ 发散指数”的指数型尺度。

\subsection{常用主部}

当 $h\to0$ 时：
\[
\begin{aligned}
\sin h&\sim h,& \tan h&\sim h,&
\arcsin h&\sim h,& \arctan h&\sim h,\\
e^h-1&\sim h,& \ln(1+h)&\sim h,&
\ln\!\left(h+\sqrt{1+h^2}\right)&\sim h,& a^h-1&\sim h\ln a,\\
{(1+h)}^\alpha-1&\sim \alpha h,&
1-\cos h&\sim \frac{h^2}{2},&
1-(\cos h)^\alpha&\sim \frac{\alpha}{2}h^2,&
\sqrt[m]{1+h}-1&\sim \frac{h}{m}.&&
\end{aligned}
\]
其中 $a>0$、$a\ne1$，$\alpha\ne0$ 为固定常数，$m\in\mathbb N_+$；所有幂、根式与对数仍必须满足相应实数定义域资格。
更一般地，若 $u\to1$，则
\[
\ln u\sim u-1.
\]
若 $\alpha(x)\to0$、$\alpha(x)\beta(x)\to0$，并且 $1+\alpha(x)$ 最终为正、$\alpha(x)\beta(x)$ 在相应去心邻域内非零，则
\[
\bigl[1+\alpha(x)\bigr]^{\beta(x)}-1\sim \alpha(x)\beta(x).
\]
它来自先取对数：$\beta\ln(1+\alpha)\sim\alpha\beta$，再用 $e^t-1\sim t$；因此这条结论的资格不是可省略的装饰。

主部发生加减抵消时，常见下一阶锚点为
\[
h-\sin h\sim\frac{h^3}{6},\qquad
\tan h-h\sim\frac{h^3}{3},\qquad
\arcsin h-h\sim\frac{h^3}{6},\qquad
h-\arctan h\sim\frac{h^3}{3}.
\]
这些高阶主部的生成与余项控制调用专题03的有限泰勒模型；本册只拥有它们作为尺度比较结果的语义。

\subsection{无穷远增长阶}

当 $x\to+\infty$，$k,\alpha>0$、$b>1$ 时，
\[
{(\ln x)}^k\ll x^\alpha\ll b^x,
\qquad f\ll g\ \text{表示}\ \frac{f}{g}\to0.
\]

因此多项式压过对数，指数压过任意固定次幂。这个序列只比较主导尺度；若表达式含正负抵消、振荡或定义域分支，仍须返回正文的存在性与尺度账本。

数列题还常用更长的增长阶梯。固定 $m,p>0$、$a>1$ 时，
\[
(\ln n)^m\ll n^p\ll a^n\ll n!\ll n^n.
\]
它同样只比较主导增长速度；阶乘可优先通过相邻项比值理解，$n^n$ 可通过取对数转成 $n\ln n$ 理解，不能把增长阶梯替代具体的符号、抵消和存在性审计。

\noindent\textbf{来源说明。} 本册由归档笔记 I-01《数列计算》、I-02《函数极限理论与计算方法》、I-03《函数连续》做来源差异与归属差异核对后重组。详细迁移台账保留在 Evidence 层；正文状态仍待使用者审阅，做题动作仍待陌生题验证。

\end{document}
